\setcounter{page}{1}
\setcounter{figure}{0}
\setcounter{chapter}{1}
%\setcounter{table}{0}



\chapter{Operaciones vectoriales}
\vspace{-14ex}
%\resp{.7\textwidth}{Fecha de entrega: \textsf{introducir fecha aqu\'i}}
\begin{flushright}
\linea{2}
\end{flushright}
\raggedright
\begin{enumerate}
\item Considere los  vectores:
$$\VEC{A}=(2\units{m})\HAT{\iota}+(-5\units{m})\HAT{\jmath}-(1\units{m})\HAT{\kappa} \qquad \mbox{y} \qquad\VEC{B}=(6\units{m}, -2\units{m}, 0\units{m}). $$

\begin{enumerate}
\item representa graficamente ambos vectores,
\item determine el \'angulo que se forma entre los dos vectores,
%\item calcule el \'area del paralelogramo formado por estos vectores.
\end{enumerate}

\item  Dibuje y exprese en notaci\'on polar (magnitud + direcci\'on) cada uno de los vectores especificados por las siguientes componentes cartesianas:
\begin{enumerate}
\item $A_x=-8.60 \, \mbox{cm}$, $A_y=-5.20 \, \mbox{cm}$.
\item $B_x=-9.70 \, \mbox{m}$, $B_y=2.45 \, \mbox{m}$.
\item $C_x=7.75 \, \mbox{km}$, $C_y=-2.70 \, \mbox{km}$.
\end{enumerate}

%\item Una fuerza $\VEC{F}=(50\units{N}; 30^{\circ})$ act\'ua a lo largo de un desplazamiento $\VEC{s}=(4.0\units{m})\HAT{\iota}-(3.0\units{N})\HAT{\jmath}$. Determine el trabajo, $W$ que efect\'ua dicha fuerza, calculando $W=\VEC{F}\cdot\VEC{s}$.
%\begin{enumerate}
%\item $\VEC{A}\cdot\VEC{B}$,
%\item $\VEC{A}\times \VEC{B}$,
%\item $|\VEC{A}\times \VEC{B}|$.
%\end{enumerate}
%\item  Un dr\'on realiza el desplazamiento  $$\Delta \VEC{r}=(1.5\U{km})\HAT{\iota}+(3.0\U{km})\HAT{\jmath}+(0.5\U{km})\HAT{\kappa}.$$ 
%\begin{enumerate}
%\item Determine un vector unitario en la direcci\'on del desplazamiento.
%\item Calcule un vector perpendicular al plano que forma $\Delta \VEC{r}$ con el eje $y$,
%\item Calcule los \textit{cosenos directores} de $\Delta \VEC{r}$.
%\end{enumerate}

%\item  En una prueba biom\'etrica, un joven corre 150 m hacia el NO, 325 m 15$^{\circ}$ SE y finalmente regresa a su punto de salida. Determine la magnitud y direcci\'on del \'ultimo desplazamiento. Exprese el resultado gr\'afica y anal\'iticamente.
\item  \label{item:vectores} En una carrera, una joven corre 1.5 km en direcci\'on 35$^{\circ}$SE, luego 1.3 km en direcci\'on 15$^{\circ}$ NO. Finalmente realiza un tercer desplazamiento, terminando a  2.4 km al NE de su punto de partida. \\ [.5cm]

\begin{enumerate}
    \item \label{item:vectores-a} Describa gr\'aficamente los dos primeros desplazamientos. %\textcolor{gray}{[utilice un sistema de referencia cartesiano]}
    \item \label{item:vectores-b}  Escriba vectorialmente los dos primeros desplazamientos. %\textcolor{gray}{[calcule las componentes rectangulares de cada vector]}
    \item \label{item:vectores-c} Determine vectorialmente, el tercer desplazamiento. %\textcolor{gray}{[calcule la suma vectorial de los tres desplazamientos]}
\end{enumerate}

 %\item \label{item:3-1}Un encargado postal realiza el recorrido que se muestra en la Figura \ref{fig:postal}. Determine la magnitud y direcci\'on del desplazamiento.%\footnote{desplazamiento es la diferencia entre el vector de posici\'on final y el inicial: $\Delta \VEC{r}=\VEC{r}_{\Mbox{final}}-\VEC{r}_{\Mbox{inicial}}$} 
 %total que realiza el encargado postal.

%\begin{figure}[h!]
%\centering
%\includegraphics[width=5in]{../figuras/ruta_postal.png}
%\caption{Recorrido encargado portal. Ejercicio \ref{item:3-1}.}
%\label{fig:postal}
%\end{figure}


%\newpage
\item \label{item:3-3} Exprese los vectores que se muestran en la Figura \ref{fig:vecs2} en notaci\'on cartesiana.

\begin{figure}[ht!]
\centering
\includegraphics[width=5in]{planeamiento_didactico/figuras/vectores.pdf}
\caption{Vectores 2D. \'Item \ref{item:3-3}.}
\label{fig:vecs2}
\end{figure}
    
\item Para los vectores de la Figura \ref{fig:vecs2}, determine las distancias entre 
\begin{enumerate}
    \item $\vec{A}$ y $\vec{B}$
    \item $\vec{A}$ y $\vec{C}$
    \item $\vec{B}$ y $\vec{C}$
\end{enumerate}
\item  \label{item:3-2} Un edificio de apartamentos tiene una base cuadrada de 20.0 m de lado y una altura de 40.0 m. La base est\'a alineada en el sentido Norte-Sur (ver Figura \ref{fig:aparta}). Un inquilino se ubica en la esquina Noroeste de la azotea, mientras que otro se encuentra en la ventana de su habitaci\'on, en el centro de la cara lateral Este. Calcular

\begin{enumerate}

\item los vectores de posici\'on de cada inquilino, respecto al ``Origen'' y 
\item el desplazamiento que deber\'ia realizar el  inquilino en la ventana para llegar donde est\'a el otro en la azotea.
\end{enumerate}
\begin{figure}[ht!]
\begin{center}
\includegraphics[width=2.75in]{planeamiento_didactico/figuras/edificio.png}
\end{center}
\caption{Edificio de apartamentos. Ejercicio \ref{item:3-2}.}
\label{fig:aparta}
\end{figure}

\newpage

\item \label{item:orion} La constelaci\'on de Ori\'on es de las m\'as famosas en el cielo de Costa Rica. En la Figura \ref{fig:orion} se muestran las posiciones aparentes de las estrellas m\'as brillantes que vemos desde la Tierra. La distancia entre la Tierra y Betelgeuse es de 643 a\~nos luz y entre la Tierra y Rigel de 860 a\~nos luz. La separaci\'on angular entre Betelgeuse y Rigel es de 16$^\circ$. Utilizando m\'etodos vectoriales, calcule la distancia, en  a\~nos luz, entre Betelgeuse a Rigel. 

\begin{figure}[ht!]
%\begin{center}
\hspace{7.5ex}
\includegraphics[width=4.8in]{planeamiento_didactico/figuras/Constelacion_Orion.png}
%\end{center}
\caption{Constelaci\'on de Orion. Ejercicio \ref{item:orion}.}
\label{fig:orion}
\end{figure}

%\newpage
%\item \label{item:3-3}  Considere cuatro posiciones respecto al centro de un lago, como se muestran en la Figura \ref{fig:vecs2}. ?`a qu\'e distancia se encuentra cada ubicación de las demás?

%\begin{figure}[h!]
%\centering
%\input{\'Itemvectores-a.tex}
%\caption{Posiciones de cuatro ubicaciones respecto al centro de un lago. Ejercicio \ref{item:3-3}.}
%\label{fig:vecs2}
%\end{figure}




\includegraphics[width=5in]{planeamiento_didactico/figuras/grid.png}

 



\end{enumerate}